\documentclass[11pt]{article}
\usepackage[a4paper,margin=1in]{geometry}
\usepackage{fontspec}
\usepackage[boldfont=false]{xeCJK}
\setCJKmainfont{FandolSong-Regular.otf}
\usepackage{amsmath}
\usepackage{amssymb}
\usepackage{array}
\usepackage{booktabs}
\usepackage{graphicx}
\usepackage{adjustbox}
\usepackage{enumitem}
\setlist[itemize]{topsep=2pt,partopsep=0pt,parsep=0pt,itemsep=2pt,leftmargin=1.4em}
\setlist[enumerate]{topsep=2pt,partopsep=0pt,parsep=0pt,itemsep=2pt,leftmargin=1.6em}
\usepackage{parskip}
\usepackage{fvextra}
\fvset{breaklines=true,breakanywhere=true,fontsize=\small}
\setlength{\parindent}{0pt}
\usepackage{longtable}
\setlength\LTpre{0pt}
\setlength\LTpost{\baselineskip}

\begin{document}

\begin{center}
  {\LARGE\bfseries Coordinate geometry}\\[0.35em]
  {\large Igcse Mathematics · Vocabulary}
\end{center}
\vspace{0.6em}

\begin{longtable}{@{}p{0.40\linewidth}p{0.24\linewidth}p{0.30\linewidth}@{}}
\toprule
\textbf{English} & \textbf{中文} & \textbf{Pinyin} \\
\midrule
\endhead
\bottomrule
\endlastfoot
coordinates & 坐标 & zuò biāo \\
axes & 坐标轴 & zuò biāo zhóu \\
horizontal & 水平 & shuǐ píng \\
vertical & 竖直 & shù zhí \\
origin & 原点 & yuán diǎn \\
quadrants & 象限 & xiàng xiàn \\
gradient & 斜率 & xié lǜ \\
intercept & 截距 & jié jù \\
table of values & 数值表 & shù zhí biǎo \\
line segment & 线段 & xiàn duàn \\
length & 长度 & cháng dù \\
Pythagoras' theorem & 勾股定理 & gōu gǔ dìng lǐ \\
midpoint & 中点 & zhōng diǎn \\
parallel & 平行 & píng xíng \\
perpendicular & 垂直 & chuí zhí \\
right angle & 直角 & zhí jiǎo \\
perpendicular bisector & 垂直平分线 & chuí zhí píng fēn xiàn \\
\end{longtable}

\end{document}
